\noindent\textbf{Por qué se usa este procedimiento.} La torsión circular cumple $\tau=Tr/J$. Para sección maciza $J=\pi D^4/32$ y $r_{max}=D/2$; sustituir y cancelar un factor $D$ da $\tau_{max}=16T/(\pi D^3)$.

\begin{enumerate}

\item $J = \pi D^4/32$ y $R = D/2$, de modo que $\tau_{max} = 16 M_x/(\pi D^3)$.

\end{enumerate}

\noindent\textbf{Operaciones y resultados intermedios.}

\begin{enumerate}

\item Tensión tangencial máxima.

\[\fitmath{y_{1}=\frac{16\cdot 500000}{\pi\cdot {40}^{3}}\approx 39.78874\,\mathrm{MPa}}\]

\end{enumerate}
